\documentclass[11pt]{article}
% Shared preamble for the printable mini lessons (1.4 onward).
%   \documentclass[11pt]{article}
%   % Shared preamble for the printable mini lessons (1.4 onward).
%   \documentclass[11pt]{article}
%   % Shared preamble for the printable mini lessons (1.4 onward).
%   \documentclass[11pt]{article}
%   \input{../preamble}
% Build the key with:  pdflatex -jobname=<name>_KEY "\def\iskey{}\input{<name>}"
% Figures shared with the web page live in figures/Unit_X/<name>.tex and are
% pulled in with \fig{Unit_X/<name>}; the standalone package drops their own
% preamble, so their macros come from figures/preamble.tex, loaded here.
\input{../../figures/preamble}
\usepackage[margin=0.75in]{geometry}
\usepackage{array}
\usepackage{enumitem}
\usepackage{amsmath}
\usepackage[helvet]{sfmath}  % math in the same sans font as the text
\usepackage{standalone}
\pagestyle{empty}
\setlength{\parindent}{0pt}
\setlist[enumerate,1]{leftmargin=1.5em,itemsep=1em}
\setlist[enumerate,2]{label=(\alph*),leftmargin=2em,itemsep=2.2em}

\newif\ifkey
\ifdefined\iskey \keytrue \fi
\colorlet{keyred}{red!75!black}
\newcommand{\ans}[1]{\ifkey\leavevmode{\color{keyred}#1}\fi}
% vertical work space on the student copy, less on the key
\newcommand{\work}[1]{\ifkey\vspace{0.6em}\else\vspace{#1}\fi}
\newcommand{\fig}[1]{\input{../../figures/#1}}
% Title line: \lessontitle{Quartiles, IQR, and Box Plots}
\newcommand{\lessontitle}[1]{{\Large\bfseries #1\ifkey\ \textcolor{keyred}{KEY}\fi}\hfill Name \underline{\hspace{2in}}\par\vspace{0.6em}}
\newcommand{\partner}{\par\vspace{0.4em}Work with a partner. Compare answers when you finish each problem.\par}

% Build the key with:  pdflatex -jobname=<name>_KEY "\def\iskey{}\input{<name>}"
% Figures shared with the web page live in figures/Unit_X/<name>.tex and are
% pulled in with \fig{Unit_X/<name>}; the standalone package drops their own
% preamble, so their macros come from figures/preamble.tex, loaded here.
% Shared preamble for web figures.  Each figure is a standalone .tex that
% \input's this file; scripts/make-figures.sh turns it into an SVG.
\usepackage[T1]{fontenc}
\usepackage{helvet}
\renewcommand{\familydefault}{\sfdefault}
\usepackage{tikz}
\usetikzlibrary{decorations.pathreplacing}
% axis numbers in the text font, not math italic
\pgfkeys{/pgf/number format/.cd,assume math mode=true,1000 sep={}}

% dot plot: \dotplot{min}{max}{value/count,...}{axis label}
\newcommand{\dotplot}[4]{%
\begin{tikzpicture}[x=0.9cm,y=0.33cm]
  \draw[thick] (#1-0.5,0) -- (#2+0.5,0);
  \foreach \t in {#1,...,#2} {\draw (\t,0.25) -- (\t,-0.25) node[below] {\t};}
  \foreach \v/\c in {#3} {\ifnum\c>0 \foreach \k in {1,...,\c} {\fill (\v,\k*1.15) circle (0.16cm);}\fi}
  \node at ({(#1+#2)/2},-2.3) {#4};
\end{tikzpicture}}

% shape sketch histogram (12 bins): \shapehist{heights}{label}
\newcommand{\shapehist}[2]{%
  \foreach \h [count=\i from 0] in {#1} \draw[fill=gray!30] (\i,0) rectangle (\i+1,\h);
  \draw (-0.3,0) -- (12.3,0);
  \node[font=\bfseries] at (6,-2.6) {#2};}
\newcommand{\meanmark}[1]{\fill (#1,0) -- ++(-0.3,-0.75) -- ++(0.6,0) -- cycle;}
\newcommand{\medmark}[2]{\draw[very thick,dashed] (#1,0) -- (#1,#2);}

% number line with minor ticks: \nline{min}{max}{label step}{axis label}
\newcommand{\nline}[4]{%
  \draw[thick] (#1-0.5,0) -- (#2+0.5,0);
  \foreach \t in {#1,...,#2} \draw (\t,0.12) -- (\t,-0.12);
  \pgfmathtruncatemacro{\nxt}{#1+#3}
  \foreach \t in {#1,\nxt,...,#2} \draw (\t,0.2) -- (\t,-0.2) node[below] {\t};
  \node at ({(#1+#2)/2},-0.95) {#4};}
% box plot above the line: \bplot{min}{Q1}{med}{Q3}{max}
\newcommand{\bplot}[5]{%
  \draw[thick] (#1,1.0) -- (#2,1.0);
  \draw[thick] (#4,1.0) -- (#5,1.0);
  \draw[thick] (#1,0.75) -- (#1,1.25);
  \draw[thick] (#5,0.75) -- (#5,1.25);
  \draw[thick,fill=gray!20] (#2,0.55) rectangle (#4,1.45);
  \draw[very thick] (#3,0.55) -- (#3,1.45);}

% ---- added for 1.4 onward ----
% axis on y=0: \axis{min}{max}{label step}{tick step}{axis label}
\newcommand{\axis}[5]{%
  \draw[thick] (#1,0) -- (#2,0);
  \pgfmathsetmacro{\axT}{#1+#4}\foreach \t in {#1,\axT,...,#2} \draw (\t,0.1) -- (\t,-0.1);
  \pgfmathsetmacro{\axL}{#1+#3}\foreach \t in {#1,\axL,...,#2}
    \draw (\t,0.18) -- (\t,-0.18) node[below] {\pgfmathprintnumber{\t}};
  \node at ({(#1+#2)/2},-0.95) {#5};}
% box plot centred at height y: \bplotat{y}{min}{Q1}{med}{Q3}{max}
\newcommand{\bplotat}[6]{%
  \draw[thick] (#2,#1) -- (#3,#1);
  \draw[thick] (#5,#1) -- (#6,#1);
  \draw[thick] (#2,#1-0.22) -- (#2,#1+0.22);
  \draw[thick] (#6,#1-0.22) -- (#6,#1+0.22);
  \draw[thick,fill=gray!20] (#3,#1-0.42) rectangle (#5,#1+0.42);
  \draw[very thick] (#4,#1-0.42) -- (#4,#1+0.42);}
% an outlier dot: \outlier{x}{y}
\newcommand{\outlier}[2]{\fill (#1,#2) circle (2.4pt);}

\usepackage[margin=0.75in]{geometry}
\usepackage{array}
\usepackage{enumitem}
\usepackage{amsmath}
\usepackage[helvet]{sfmath}  % math in the same sans font as the text
\usepackage{standalone}
\pagestyle{empty}
\setlength{\parindent}{0pt}
\setlist[enumerate,1]{leftmargin=1.5em,itemsep=1em}
\setlist[enumerate,2]{label=(\alph*),leftmargin=2em,itemsep=2.2em}

\newif\ifkey
\ifdefined\iskey \keytrue \fi
\colorlet{keyred}{red!75!black}
\newcommand{\ans}[1]{\ifkey\leavevmode{\color{keyred}#1}\fi}
% vertical work space on the student copy, less on the key
\newcommand{\work}[1]{\ifkey\vspace{0.6em}\else\vspace{#1}\fi}
\newcommand{\fig}[1]{\input{../../figures/#1}}
% Title line: \lessontitle{Quartiles, IQR, and Box Plots}
\newcommand{\lessontitle}[1]{{\Large\bfseries #1\ifkey\ \textcolor{keyred}{KEY}\fi}\hfill Name \underline{\hspace{2in}}\par\vspace{0.6em}}
\newcommand{\partner}{\par\vspace{0.4em}Work with a partner. Compare answers when you finish each problem.\par}

% Build the key with:  pdflatex -jobname=<name>_KEY "\def\iskey{}\input{<name>}"
% Figures shared with the web page live in figures/Unit_X/<name>.tex and are
% pulled in with \fig{Unit_X/<name>}; the standalone package drops their own
% preamble, so their macros come from figures/preamble.tex, loaded here.
% Shared preamble for web figures.  Each figure is a standalone .tex that
% \input's this file; scripts/make-figures.sh turns it into an SVG.
\usepackage[T1]{fontenc}
\usepackage{helvet}
\renewcommand{\familydefault}{\sfdefault}
\usepackage{tikz}
\usetikzlibrary{decorations.pathreplacing}
% axis numbers in the text font, not math italic
\pgfkeys{/pgf/number format/.cd,assume math mode=true,1000 sep={}}

% dot plot: \dotplot{min}{max}{value/count,...}{axis label}
\newcommand{\dotplot}[4]{%
\begin{tikzpicture}[x=0.9cm,y=0.33cm]
  \draw[thick] (#1-0.5,0) -- (#2+0.5,0);
  \foreach \t in {#1,...,#2} {\draw (\t,0.25) -- (\t,-0.25) node[below] {\t};}
  \foreach \v/\c in {#3} {\ifnum\c>0 \foreach \k in {1,...,\c} {\fill (\v,\k*1.15) circle (0.16cm);}\fi}
  \node at ({(#1+#2)/2},-2.3) {#4};
\end{tikzpicture}}

% shape sketch histogram (12 bins): \shapehist{heights}{label}
\newcommand{\shapehist}[2]{%
  \foreach \h [count=\i from 0] in {#1} \draw[fill=gray!30] (\i,0) rectangle (\i+1,\h);
  \draw (-0.3,0) -- (12.3,0);
  \node[font=\bfseries] at (6,-2.6) {#2};}
\newcommand{\meanmark}[1]{\fill (#1,0) -- ++(-0.3,-0.75) -- ++(0.6,0) -- cycle;}
\newcommand{\medmark}[2]{\draw[very thick,dashed] (#1,0) -- (#1,#2);}

% number line with minor ticks: \nline{min}{max}{label step}{axis label}
\newcommand{\nline}[4]{%
  \draw[thick] (#1-0.5,0) -- (#2+0.5,0);
  \foreach \t in {#1,...,#2} \draw (\t,0.12) -- (\t,-0.12);
  \pgfmathtruncatemacro{\nxt}{#1+#3}
  \foreach \t in {#1,\nxt,...,#2} \draw (\t,0.2) -- (\t,-0.2) node[below] {\t};
  \node at ({(#1+#2)/2},-0.95) {#4};}
% box plot above the line: \bplot{min}{Q1}{med}{Q3}{max}
\newcommand{\bplot}[5]{%
  \draw[thick] (#1,1.0) -- (#2,1.0);
  \draw[thick] (#4,1.0) -- (#5,1.0);
  \draw[thick] (#1,0.75) -- (#1,1.25);
  \draw[thick] (#5,0.75) -- (#5,1.25);
  \draw[thick,fill=gray!20] (#2,0.55) rectangle (#4,1.45);
  \draw[very thick] (#3,0.55) -- (#3,1.45);}

% ---- added for 1.4 onward ----
% axis on y=0: \axis{min}{max}{label step}{tick step}{axis label}
\newcommand{\axis}[5]{%
  \draw[thick] (#1,0) -- (#2,0);
  \pgfmathsetmacro{\axT}{#1+#4}\foreach \t in {#1,\axT,...,#2} \draw (\t,0.1) -- (\t,-0.1);
  \pgfmathsetmacro{\axL}{#1+#3}\foreach \t in {#1,\axL,...,#2}
    \draw (\t,0.18) -- (\t,-0.18) node[below] {\pgfmathprintnumber{\t}};
  \node at ({(#1+#2)/2},-0.95) {#5};}
% box plot centred at height y: \bplotat{y}{min}{Q1}{med}{Q3}{max}
\newcommand{\bplotat}[6]{%
  \draw[thick] (#2,#1) -- (#3,#1);
  \draw[thick] (#5,#1) -- (#6,#1);
  \draw[thick] (#2,#1-0.22) -- (#2,#1+0.22);
  \draw[thick] (#6,#1-0.22) -- (#6,#1+0.22);
  \draw[thick,fill=gray!20] (#3,#1-0.42) rectangle (#5,#1+0.42);
  \draw[very thick] (#4,#1-0.42) -- (#4,#1+0.42);}
% an outlier dot: \outlier{x}{y}
\newcommand{\outlier}[2]{\fill (#1,#2) circle (2.4pt);}

\usepackage[margin=0.75in]{geometry}
\usepackage{array}
\usepackage{enumitem}
\usepackage{amsmath}
\usepackage[helvet]{sfmath}  % math in the same sans font as the text
\usepackage{standalone}
\pagestyle{empty}
\setlength{\parindent}{0pt}
\setlist[enumerate,1]{leftmargin=1.5em,itemsep=1em}
\setlist[enumerate,2]{label=(\alph*),leftmargin=2em,itemsep=2.2em}

\newif\ifkey
\ifdefined\iskey \keytrue \fi
\colorlet{keyred}{red!75!black}
\newcommand{\ans}[1]{\ifkey\leavevmode{\color{keyred}#1}\fi}
% vertical work space on the student copy, less on the key
\newcommand{\work}[1]{\ifkey\vspace{0.6em}\else\vspace{#1}\fi}
\newcommand{\fig}[1]{\input{../../figures/#1}}
% Title line: \lessontitle{Quartiles, IQR, and Box Plots}
\newcommand{\lessontitle}[1]{{\Large\bfseries #1\ifkey\ \textcolor{keyred}{KEY}\fi}\hfill Name \underline{\hspace{2in}}\par\vspace{0.6em}}
\newcommand{\partner}{\par\vspace{0.4em}Work with a partner. Compare answers when you finish each problem.\par}

% Unit summary: the unit's tools in tables, to print and keep. Same tables as
% chapters/Unit_7/Summary.qmd; edit both together. Nothing to answer, so the
% _KEY copy is the same sheet.
\newcommand{\sumhead}[1]{\par\vspace{0.7em}{\large\bfseries #1}\par\vspace{0.3em}}
\newcolumntype{L}[1]{>{\raggedright\arraybackslash}p{#1}}
\renewcommand{\arraystretch}{1.25}
% a small distribution table; [1.5] gives fractions room to clear the rules
\newcommand{\disttable}[3][1.25]{{\renewcommand{\arraystretch}{#1}\begin{tabular}{#2}#3\end{tabular}}}

\begin{document}

\lessontitle{Unit 7 Summary: Probability Distributions}

\sumhead{Random Variables and Distributions}
Flip a coin 2 times. Let $X$ = the number of heads. \qquad
\disttable[1.5]{|c|c|c|c|}{\hline
$x$ (heads) & 0 & 1 & 2\\ \hline
$P(x)$ & $\frac{1}{4}$ & $\frac{2}{4} = \frac{1}{2}$ & $\frac{1}{4}$\\ \hline}

\vspace{0.5em}
\begin{tabular}{|L{1.45in}|L{2.6in}|L{2.6in}|}\hline
\textbf{Idea} & \textbf{What it means} & \textbf{Example: flip a coin 2 times}\\ \hline
Random variable & A number that comes from a chance process. & $X$ = the number of heads.\\ \hline
$X$, $x$, and $P(x)$ & Capital $X$ is the random variable. Small $x$ is one of its values. $P(x)$ means the probability that $X$ equals $x$. & $x$ is 0, 1, or 2. \; $P(1) = \frac{1}{2}$.\\ \hline
Probability distribution & It lists every value of $X$ and its probability. It can be a table or a graph. & The table above.\\ \hline
Valid distribution & 1. Every probability is from 0 to 1.\par 2. The probabilities add to 1: $\Sigma P(x) = 1$ ($\Sigma$ means add up). & $\frac{1}{4} + \frac{1}{2} + \frac{1}{4} = 1$. Valid.\\ \hline
Using a distribution & Add the probabilities of the values you want. & $P(X \ge 1) = P(1) + P(2)$\par $= \frac{1}{2} + \frac{1}{4} = \frac{3}{4}$\\[0.2em] \hline
\end{tabular}

\vspace{0.3em}
\textbf{Symbols:} $\ge$ means ``greater than or equal to.'' \quad $\le$ means ``less than or equal to.''

\sumhead{Theoretical or Empirical?}
\begin{tabular}{|L{1.45in}|L{2.6in}|L{2.6in}|}\hline
\textbf{Kind} & \textbf{Where it comes from} & \textbf{Example}\\ \hline
Theoretical & Reasoning about the chance process (equally likely outcomes, or the multiplication rule). & The coin flips above.\\ \hline
Empirical & Data. Divide each count by the total. & 50 families: 0 cars: 5, 1 car: 20, 2 cars: 18, 3 cars: 7. So $P(0) = \frac{5}{50} = 0.10$, $P(1) = 0.40$, $P(2) = 0.36$, $P(3) = 0.14$.\\[0.2em] \hline
\end{tabular}

\sumhead{Expected Value}
The number of goals a soccer team scores in a game: \qquad
\disttable{|c|c|c|c|c|c|}{\hline
$x$ (goals) & 0 & 1 & 2 & 3 & 4\\ \hline
$P(x)$ & 0.1 & 0.3 & 0.3 & 0.2 & 0.1\\ \hline
$x \cdot P(x)$ & 0 & 0.3 & 0.6 & 0.6 & 0.4\\ \hline}

\vspace{0.5em}
\begin{tabular}{|L{1.45in}|L{2.6in}|L{2.6in}|}\hline
\textbf{Tool} & \textbf{How it works} & \textbf{Example: goals in a game}\\ \hline
Expected value & $E(X) = \Sigma\, x \cdot P(x)$. \; $E(X)$ means the expected value of $X$. Multiply each value by its probability. Then add them all up. & $E(X) = 0 + 0.3 + 0.6 + 0.6 + 0.4$\par $= 1.9$ goals\\ \hline
Mean of $X$ & The expected value is also called the \textbf{mean} of $X$, written $\mu_X$ ($\mu$ is ``mu''). & $\mu_X = 1.9$ goals\\ \hline
What it means & A \textbf{long-run average}. It tells you what happens over many tries. It does not tell you what happens on one try. & Over many games, the team scores about 1.9 goals per game on average. No single game has 1.9 goals.\\ \hline
\end{tabular}

\newpage
\sumhead{Fair or Not?}
\begin{tabular}{|L{1.45in}|L{2.6in}|L{2.6in}|}\hline
 & \textbf{What it means} & \textbf{Example}\\ \hline
Fair method of choosing & Every person has the same probability of being chosen. & Roll a number cube. 1 or 2: Ana, 3 or 4: Ben, 5 or 6: Cai. Each has $\frac{2}{6} = \frac{1}{3}$. \textbf{Fair.}\\[0.2em] \hline
Not a fair method & Not every person has the same probability. & Flip 2 coins. HH: Ana, TT: Ben, one of each: Cai. Ana $\frac{1}{4}$, Ben $\frac{1}{4}$, Cai $\frac{2}{4} = \frac{1}{2}$. \textbf{Not fair.}\\[0.2em] \hline
Fair game & The expected value of the money you win or lose is \textbf{0}. On average, no one wins or loses money. & Roll a number cube. You win \$3 if you roll a 6. You pay \$1 for any other number.\par $E(X) = 3 \cdot \frac{1}{6} + (-1) \cdot \frac{5}{6} = -\frac{2}{6} \approx -\$0.33$\par You lose about 33 cents a game. \textbf{Not fair.}\\[0.2em] \hline
\end{tabular}

\sumhead{Expected Payoff}
A lottery ticket costs \$2. 1 in 1,000 tickets wins \$500. 1 in 50 tickets wins \$10. The rest win nothing.

\vspace{0.4em}
\disttable{|l|c|c|c|}{\hline
Outcome & Net gain $x$ & $P(x)$ & $x \cdot P(x)$\\ \hline
Win \$500 & $500 - 2 = 498$ & 0.001 & 0.498\\ \hline
Win \$10 & $10 - 2 = 8$ & 0.02 & 0.16\\ \hline
Win nothing & $-2$ & 0.979 & $-1.958$\\ \hline
\multicolumn{3}{|r|}{$E(X) = \Sigma\, x \cdot P(x)$} & $-1.30$\\ \hline}

\vspace{0.5em}
\begin{tabular}{|L{1.45in}|L{2.6in}|L{2.6in}|}\hline
\textbf{Tool} & \textbf{How it works} & \textbf{Example: the lottery ticket}\\ \hline
Net gain & Net gain $=$ what you win $-$ what you paid. It is negative when you lose money. & Win \$500: $500 - 2 = 498$.\par Win nothing: $-2$.\\ \hline
Expected payoff & Let $X$ = the net gain. Then $E(X) = \Sigma\, x \cdot P(x)$. & $0.498 + 0.16 + (-1.958) = -1.30$\\ \hline
What it means & The amount you win or lose on average. & On average, a buyer \textbf{loses \$1.30} per ticket. The lottery keeps that money.\\ \hline
\end{tabular}

\vspace{0.3em}
\textbf{Negative expected payoff:} you lose money on average. \quad \textbf{Expected payoff of 0:} the game is fair.

\sumhead{Protection Plans}
\begin{tabular}{|L{1.45in}|L{2.6in}|L{2.6in}|}\hline
\textbf{Step} & \textbf{What to do} & \textbf{Example: a phone}\\ \hline
Expected cost with no plan & Multiply the chance it breaks by the cost to replace it. & A phone costs \$300 to replace. The chance it breaks this year is 0.10.\par $0.10 \cdot 300 = \$30$\\ \hline
Compare & Compare it to the cost of the plan. & The plan costs \$40 for the year. On average, the plan costs \$10 more. That is how the company makes money.\\ \hline
Think about risk & Expected value is one part of a decision. Risk is another part: could you pay the big bill if it came? & The plan protects you from a sudden \$300 bill. If you could not pay it, the plan might be worth it to you.\\ \hline
\end{tabular}

\end{document}
